\begin{table*}[t]
\centering
\caption{Supplementary Table S1: Comparative Overview of Representative Threat Detection Studies vs. Proposed ScamShield-VN}
\label{tab:supp_related_work}
\resizebox{\textwidth}{!}{%
\begin{tabular}{lllllp{5.0cm}}
\toprule
\textbf{Study} & \textbf{Language} & \textbf{Model / Architecture} & \textbf{Evaluation Dataset} & \textbf{Key Metric} & \textbf{Primary Limitation / Research GAP} \\
\midrule
Tuan et al. (2023) & Vietnamese & PhoBERT, CNN, LSTM & Multi-accent SMS (N~3,000) & Acc: 95.56\% & Evaluated under symmetric loss; vulnerable to noisy word-segmentation failures. \\
Nguyen-Xuan et al. (2026) & Viet-Eng & MLTEA Dual-Encoder & Code-mixed Phishing & F1: 0.940 & Cloud-scale dual-encoder; prohibitive inference latency for edge mobile integration. \\
Cam & Binh (2026) & Vietnamese & PhoBERT, BiLSTM & VNSED Email Corpus & Acc: 86.47\% & Limited coverage of colloquial teencode; elevated false negative rate on subtle spam. \\
Zhou et al. (2026) & Multilingual & GPT-4o, Claude, Gemini & FraudSMSWalker (N~4,000) & Recall: 64.16--90.13\% & Severe alarm fatigue on legitimate traffic (benign recall collapsed to 12.81--30.25\%). \\
Banda et al. (2026) & English & Distilled SLM (Llama-3.2-3B) & QR/Text Lures (N=4,200) & Latency: 25.3--45 ms & Focused on multi-modal English lures; non-trivial mobile memory footprint. \\
Phap et al. (2026) & English & BiGRU + MaskedPool + WBCE & UCI & LSDST (N=5,574) & Latency: 0.25 ms & Recurrent architecture evaluated on English; lacks self-attention context for low-resource slang. \\
This Work (ScamShield-VN) & Vietnamese & Two-Tier: ViSoBERT + 3-Agent Council & Unified SMS (N=2,664 / 384) & Recall: 97.11\%, FPR: 2.37\% & Exploratory trade-off: 267.5 ms, 0.00\% bank FPR (N=19), \$0.0153/1k SMS, edge privacy preservation. \\
\bottomrule
\end{tabular}%
}
\par\vspace{6pt}
\parbox{\textwidth}{\footnotesize \textit{Note:} Reported metrics are extracted from corresponding publications evaluated on heterogeneous datasets, hardware environments, and task formulations; they illustrate comparative architectural paradigms rather than direct benchmark parity.}
\end{table*}
